\begin{description}
	\item[1行目] 乱数を10万個作ることにします。
	\item[2行目] 標準正規分布の乱数発生
	\item[3行目] データフレーム型に整形
	\item[4行目] 変数名がつけられていないので，\texttt{val}という名前をつけることに。ここでの$1$は1列目を意味している。
	\item[6行目から] データフレーム\texttt{X}を使って計算
	\item[7行目] \texttt{summarise}関数は要約計算。変数名$=$計算式で表す。
	\item[8行目] 期待値を\texttt{mean}関数で算出
	\item[9行目] 確率分布の標準偏差を\texttt{sd}関数で算出\footnote{Rの\texttt{sd}関数は不偏分散の平方根を計算しており，厳密には標本標準偏差の計算で良いが，サンプルサイズが十分に大きいのでほとんど影響しません。}
	\item[10行目] 中央値を\texttt{median}関数で計算
	\item[11-13行目] パーセンタイル点を算出。50パーセンタイル点は中央値に同じ。
\end{description}
